\section{Ablation studies}
\label{sec:ablation}

\textbf{Effect of the expert-cost weight: }
Figure~\ref{fig:alpha_sweep} shows that increasing $\alpha$ reduces both
the mean selected cardinality and mIoU. Our method generally achieves
higher mIoU than the baselines at comparable correction budgets,
showing a consistent quality--cost trade-off across datasets.

\noindent\textbf{Number of sampled frame subsets: }
Table~\ref{tab:number_random_subsets} shows that increasing the number of
sampled subsets generally improves mIoU. Compared with no sampled alternatives
(\(B=0\)), \(B=4\) improves mIoU by 0.034 on VOST, 0.022 on CAMotion, and
0.032 on VTUS. SUN-SEG performs best at \(B=2\), with little change at \(B=4\).
We therefore use \(B=4\) in the main experiments.

\noindent\textbf{Cardinality and frame-selection components:}
Table~\ref{tab:component_ablation} shows that replacing either learned
component with a random alternative reduces video-level mIoU. Randomising
the cardinality decision causes the larger reduction, particularly on
VOST and VTUS, while random frame selection also lowers performance
across all datasets. Combining the learned cardinality model with
multi-frame selection achieves the highest mIoU on every dataset.

\noindent\textbf{Multiple correction rounds:}
Although our framework requests all corrections in a single round, the same policy can be reapplied after each re-propagation.
Figure~\ref{fig:sunseg_iterative_prompting} reports this setting on SUN-SEG for maximum cardinalities \(K\in\{1,2,3,4\}\).
Larger batches reach a given mean IoU in fewer interaction rounds, which matters because each round returns the video to the expert for review.
Each round applies the policy to the current propagation alone, without distributing a correction budget across rounds, and optimising jointly across rounds is a natural extension.
